\Chapter{54}

\Verse{1}
{
    \zA{却说贾珍贾琏暗暗预备下大笸箩的钱，听见贾母说赏，忙命小厮们快撒钱，只 听满台钱响，贾母大悦。 二人遂起身，小厮们忙将一把新暖银壶捧来，递与贾琏手
        内，随了贾珍趋至里面。 贾珍先到李婶娘席上，躬身取下杯来，回身，贾琏忙斟了 一盏，然后便至薛姨妈席上也斟了。 二人忙起来笑说：“二位爷请坐着罢了，何必
        多礼。” 于是除邢王二夫人，满席都离了席，也俱垂手旁站。 贾珍等至贾母榻前， 因榻矮，二人便屈膝跪了，贾珍在前捧杯，贾琏在后捧壶。}
}
{
    \eA{Chia Chen and Chia Lien had, we will now explain, secretly got ready large
        baskets of cash, so the moment they heard old lady Chia utter the word
        'tip,' they promptly bade the pages be quick and fling the money. The noise
        of the cash, running on every side of the stage, was all that fell on the
        ear. Dowager lady Chia thoroughly enjoyed it. The two men then rose to their
        feet.}
}
{
}

\Verse{2}
{
    \zA{虽只二人捧酒，那贾琮 弟兄等却都是一溜排班随着他二人进来，见他二人跪下，都一溜跪下。 宝玉也忙跪 下。 湘云悄推他，笑道：“你这会子又帮着跪下做什么？
        有这么着的呢，你也去斟 一巡酒，岂不好？” 宝玉悄笑道：“再等一会再斟去。” 说着，等他二人斟完，起 来，又给邢王二夫人斟过了。
        贾珍笑说：“妹妹们怎么着呢？” 贾母等都说道：“你 们去罢，他们倒便宜些呢。” 贾珍等方退出。 当下天有二鼓，戏演的是《八义？
        观灯》八出，正在热闹之际。}
}
{
    \eA{The pages hastened to lay hold of a silver kettle, newly brought in with
        fresh wine, and to deposit it in Chia Lien's hands, who followed Chia Chen
        with quick step into the inner rooms. Chia Chen advanced first up to
        'sister-in-law' Li's table, and curtseying, he raised her cup, and turned
        round, whereupon Chia Lien quickly filled it to the brim.}
}
{
}

\Verse{3}
{
    \zA{宝玉因下席往 外走。 贾母问：“往那里去？ 外头炮仗利害，留神天上吊下火纸来烧着。” 宝玉笑 回说：“不往远去，只出去就来。”
        贾母命婆子们：“好生跟着。” 于是宝玉出来， 只有麝月秋纹几个小丫头随着。 贾母因说：“袭人怎么不见？ 他如今也有些拿大了， 单支使小女孩儿出来。”
        王夫人忙起身笑说道：“他妈前日没了，因有热孝，不便 前头来。” 贾母点头，又笑道：“跟主子，却讲不起这孝与不孝。 要是他还跟我， 难道这会子也不在这里？}
}
{
    \eA{Next they approached Mrs. Hsüeh's table, and they also replenished her cup.
        These two ladies lost no time in standing up, and smilingly expostulating.
        "Gentlemen," they said, "please take your seats. What's the use of standing
        on such ceremonies?"}
}
{
}

\Verse{4}
{
    \zA{这些竟成了例了。” 凤姐儿忙过来笑回道：“今晚便没孝， 那园子里头也须得看着灯烛花爆，最是担险的。 这里一唱戏，园子里的谁不来偷瞧
        瞧，他还细心，各处照看。 况且这一散后，宝兄弟回去睡觉，各色都是齐全的。 若 他再来了，众人又不经心，散了回去，铺盖也是冷的，茶水也不齐全，便各色都不
        便宜，自然我叫他不用来。 老祖宗要叫他来，我就叫他就是了。” 贾母听了这话， 忙说：“你这话很是，你必想的周到，快别叫他了。 但只他妈几时没了？}
}
{
    \eA{But presently every one, with the exception of the two ladies Mesdames Hsing
        and Wang, quitted the banquet and dropping their arms against their bodies
        they stood on one side. Chia Chen and his companion then drew near dowager
        lady Chia's couch. But the couch was so low that they had to stoop on their
        knees. Chia Chen was in front, and presented the cup. Chia Lien was behind,
        and held the kettle up to her.}
}
{
}

\Verse{5}
{
    \zA{我怎么不 知道？” 凤姐儿笑道：“前儿袭人去亲自回老太太的，怎么倒忘了？” 贾母想了想， 笑道：“想起来了。 我的记性竟平常了。”
        众人都笑说：“老太太那里记得这些事。” 贾母因又叹道：“我想着他从小儿伏侍我一场，又伏侍了云儿，末后给了个魔王， 给他魔了这好几年。
        他又不是咱们家根生土长的奴才，没受过咱们什么大恩典，他 娘没了，我想着要给他几两银子发送他娘，也就忘了。” 凤姐儿道：“前儿太太赏
        了他四十两银子，就是了。”}
}
{
    \eA{But notwithstanding that only these two offered her wine, Chia Tsung and the
        other young men followed them closely in the order of their age and grade;
        so the moment they saw them kneel, they immediately threw themselves on
        their knees. Pao-yü too prostrated himself at once. Hsiang-yün stealthily
        gave him a push. "What's the use of your now following their lead again and
        falling on your knees?" she said.}
}
{
}

\Verse{6}
{
    \zA{贾母听说，点头道：“这还罢了。 正好前儿鸳鸯的娘 也死了，我想他老子娘都在南边，我也没叫他家去守孝。 如今他两处全礼，何不叫 他二人一处作伴去？”
        又命婆子拿些果子菜馔点心之类与他二人吃去。 琥珀笑道： “还等这会子？ 他早就去了。” 说着，大家又吃酒看戏。
        且说宝玉一径来至园中，众婆子见他回房，便不跟去，只坐在园门里茶房里烤 火，和管茶的女人偷空饮酒斗牌。 宝玉至院中，虽是灯光灿烂，却无人声。 麝月道：
        “他们都睡了不成？}
}
{
    \eA{"But since you behave like this, wouldn't it be well if you also went and
        poured wine all round?" Pao-yü laughed. "Hold on a bit," he rejoined in a
        low tone, "and I'll go and do so." So speaking, he waited until his two
        relatives had finished pouring the wine and risen to their feet, when he
        also went and replenished the cups of Mesdames Wang and Hsing. "What about
        the young ladies?" Chia Chen smilingly asked.}
}
{
}

\Verse{7}
{
    \zA{咱们悄悄进去吓他们一跳。” 于是大家蹑手蹑脚，潜踪进镜壁 去一看，只见袭人和一个人对歪在地炕上，那一头有两个老嬷嬷打盹。 宝玉只当他
        两个睡着了，才要进去，忽听鸳鸯嗽了一声，说道：“天下事可知难定。 论理你单 身在这里，父母在外头，每年他们东去西来，没个定准，想来你是再不能送终的了；
        偏生今年就死在这里，你倒出去送了终。” 袭人道：“正是，我也想不到能够看着 父母殡殓。 回了太太，又赏了四十两银子，这倒也算养我一场，我也不敢妄想了。”}
}
{
    \eA{"You people had better be going," old lady Chia and the other ladies
        unanimously observed. "They'll, then, be more at their ease." At this hint
        Chia Chen and his companions eventually withdrew. The second watch had not,
        at the time, yet gone. The play that was being sung was: 'The eight worthies
        look at the lanterns,' consisting of eight acts; and had now reached a
        sensational part.}
}
{
}

\Verse{8}
{
    \zA{宝玉听了，忙转身悄向麝月等道：“谁知他也来了。 我这一进去，他又赌气走了， 不如咱们回去罢，让他两个清清净净的说话。 袭人正在那里闷着，幸他来的好。”
        说着，仍悄悄出来。 宝玉便走过山石后去，站着撩衣。 麝月秋纹皆站住，背过脸去， 口内笑说：“蹲下再解小衣，留神风吹了肚子。”
        后面两个小丫头知是小解，忙先 出去茶房内预备水去了。 这里宝玉刚过来，只见两个媳妇迎面来了，又问：“是谁？” 秋纹道：“宝玉
        在这里呢，大呼小叫，留神吓着罢！”}
}
{
    \eA{Pao-yü at this stage left the feast and was going out. "Where are you off
        to?" inquired his grandmother Chia. "The crackers outside are dreadful.
        Mind, the lighted pieces of paper falling from above might burn you." Pao-yü
        smiled. "I'm not going far," he answered. "I'm merely going out of the room,
        and will be back at once." Dowager lady Chia directed the matrons to "be
        careful and escort him."}
}
{
}

\Verse{9}
{
    \zA{那媳妇们忙笑道：“我们不知，大节下来惹 祸了。 姑娘们可连日辛苦了！” 说着，已到跟前。 麝月等问：“手里拿着什么？”
        媳妇道：“是老太太赏金、花二位姑娘吃的。” 秋纹笑道：“外头唱的是《八义》， 没唱《混元盒》，那里又跑出‘金花娘娘’来了？” 宝玉命：“揭起来我瞧瞧。”
        秋纹麝月忙上去将两个盒子揭开，两个媳妇忙蹲下身子。 宝玉看了两个盒内都是席 上所有的上等果品茶点，点了一点头就走。 麝月等忙胡乱掷了盒盖跟上来。}
}
{
    \eA{Pao-yü forthwith sallied out; with no other attendants however than She
        Yüeh, Ch'iu Wen and several youthful maids. "How is it," his grandmother
        Chia felt obliged so ask, "that I don't see anything of Hsi Jen? Is she too
        now putting on high and mighty airs that she only sends these juvenile girls
        here?" Madame Wang rose to her feet with all haste. "Her mother," she
        explained, "died the other day;}
}
{
}

\Verse{10}
{
    \zA{宝玉笑 道：“这两个女人倒和气，会说话。 他们天天乏了，倒说你们连日辛苦，倒不是那 矜功自伐的。” 麝月道：“这两个就好，那不知理的是太不知理。”
        宝玉道：“你 们是明白人，担待他们是粗夯可怜的人就完了。” 一面说，一面就走出了园门。 那 几个婆子虽吃酒斗牌，却不住出来打探，见宝玉出来，也都跟上来。
        到了花厅廊上， 只见那两个小丫头，一个捧着个小盆，又一个搭着手巾，又拿着沤子小壶儿，在那 里久等。}
}
{
    \eA{so being in deep mourning, she couldn't very well present herself." Dowager
        lady Chia nodded her head assentingly. "When one is in service," she
        smilingly remarked, "there should be no question of mourning or no mourning.
        Is it likely that, if she were still in my pay, she wouldn't at present be
        here? All these practices have quite become precedents!" Lady Feng crossed
        over to her.}
}
{
}

\Verse{11}
{
    \zA{秋纹先忙伸手向盆内试了试，说道：“你越大越粗心了，那里弄得这冷水？” 小丫头笑道：“姑娘瞧瞧，这个天，我怕水冷，倒的是滚水，这还冷了。” 正说着，
        可巧见一个老婆子提着一壶滚水走来，小丫头就说：“好奶奶，过来给我倒上些水。” 那婆子道：“姐姐，这是老太太沏茶的，劝你去舀罢。 那里就走大了脚呢？”
        秋纹 道：“不管你是谁的!你不给我，管把老太太的茶铞子倒了洗手！” 那婆子回头见 了秋纹，忙提起壶来倒了些。}
}
{
    \eA{"Had she even not been in mourning to-night," she chimed in with a laugh,
        "she would have had to be in the garden and keep an eye over that pile of
        lanterns, candles, and fireworks, as they're most dangerous things. For as
        soon as any theatricals are set on foot in here, who doesn't surreptitiously
        sneak out from the garden to have a look? But as far as she goes, she's
        diligent, and careful of every place.}
}
{
}

\Verse{12}
{
    \zA{秋纹道：“够了!你这么大年纪，也没见识。 谁不知 是老太太的？ 要不着的就敢要了？” 婆子笑道：“我眼花了，没认出这姑娘来。”
        宝玉洗了手，那小丫头子拿小壶儿倒了沤子在他手内，宝玉沤了。 秋纹麝月也趁热 水洗了一回，跟进宝玉来。 宝玉便要了一壶暖酒，也从李婶娘斟起。
        他二人也笑让坐。 贾母便说：“他小 人家儿，让他斟去。 大家倒要干过这杯。” 说着，便自己干了。 邢王二夫人也忙干 了，薛姨妈李婶娘也只得干了。}
}
{
    \eA{Moreover, when the company disperses and brother Pao-yü retires to sleep,
        everything will be in perfect readiness. But, had she also come, that bevy
        of servants wouldn't again have cared a straw for anything; and on his
        return, after the party, the bedding would have been cold, the tea-water
        wouldn't have been ready, and he would have had to put up with every sort of
        discomfort.}
}
{
}

\Verse{13}
{
    \zA{贾母又命宝玉道：“你连姐姐妹妹的一齐斟上，不 许乱斟，都要叫他干了。” 宝玉听说，答应着，一一按次斟上了。 至黛玉前，偏他
        不饮，拿起杯来，放在宝玉唇边。 宝玉一气饮干，黛玉笑说：“多谢。” 宝玉替他 斟上一杯。 凤姐儿便笑道：“宝玉别喝冷酒。
        仔细手颤，明儿写不的字，拉不的弓。” 宝玉道：“没有吃冷酒。” 凤姐儿笑道：“我知道没有，不过白嘱咐你。” 然后宝
        玉将里面斟完，只除贾蓉之妻是命丫鬟们斟的。 复出至廊下，又给贾珍等斟了。}
}
{
    \eA{That's why I told her that there was no need for her to come. But should
        you, dear senior, wish her here, I'll send for her straightway and have
        done." Old lady Chia lent an ear to her arguments. "What you say," she
        promptly put in, "is perfectly right. You've made better arrangements than I
        could. Quick, don't send for her! But when did her mother die? How is it I
        know nothing about it?"}
}
{
}

\Verse{14}
{
    \zA{坐 了一回，方进来，仍归旧坐。 一时上汤之后，又接着献元宵。 贾母便命：“将戏暂歇，小孩子们可怜见的， 也给他们些滚汤热菜的吃了再唱。”
        又命将各样果子元宵等物拿些给他们吃。 一时 歇了戏，便有婆子带了两个门下常走的女先儿进来，放了两张杌子在那一边，贾母 命他们坐了，将弦子琵琶递过去。
        贾母便问李薛二人：“听什么书？” 他二人都回 说：“不拘什么都好。” 贾母便问：“近来可又添些什么新书？”}
}
{
    \eA{"Some time ago," lady Feng laughed, "Hsi Jen came in person and told you,
        worthy ancestor, and how is it you've forgotten it?" "Yes," resumed dowager
        lady Chia smiling, after some reflection, "I remember now. My memory is
        really not of the best." At this, everybody gave way to laughter. "How could
        your venerable ladyship," they said, "recollect so many matters?" Dowager
        lady Chia thereupon heaved a sigh.}
}
{
}

\Verse{15}
{
    \zA{两个女先回说： “倒有一段新书，是残唐五代的故事。” 贾母问是何名，女先儿回说：“这叫做《凤 求鸾》。” 贾母道：“这个名字倒好，不知因什么起的？
        你先说大概，若好再说。” 女先儿道：“这书上乃是说残唐之时，那一位乡绅，本是金陵人氏，名唤王忠，曾
        做过两朝宰辅，如今告老还家，膝下只有一位公子，名唤王熙凤。” 众人听了，笑 将起来。 贾母笑道：“这不重了我们凤丫头了！” 媳妇忙上去推他说：“是二奶奶
        的名字，少混说。”}
}
{
    \eA{"How I remember," she added, "the way she served me ever since her youth up;
        and how she waited upon Yün Erh also; how at last she was given to that
        prince of devils, and how she has slaved away with that imp for the last few
        years. She is, besides, not a slave-girl, born or bred in the place. Nor has
        she ever received any great benefits from our hands.}
}
{
}

\Verse{16}
{
    \zA{贾母道：“你只管说罢。” 女先儿忙笑着站起来说：“我们该 死了!不知是奶奶的讳。” 凤姐儿笑道：“怕什么!你说罢。 重名重姓的多着呢。”
        女先儿又说道：“那年王老爷打发了王公子上京赶考，那日遇了大雨，到了一个庄 子上避雨。 谁知这庄上也有位乡绅，姓李，与王老爷是世交，便留下这公子住在书
        房里。 这李乡绅膝下无儿，只有一位千金小姐。 这小姐芳名叫做雏鸾，琴棋书画， 无所不通。” 贾母忙道：“怪道叫做《凤求鸾》。}
}
{
    \eA{When her mother died, I meant to have given her several taels for her
        burial; but it quite slipped from my mind." "The other day," lady Feng
        remarked, "Madame Wang presented her with forty taels; so that was all
        right." At these words, old lady Chia nodded assent. "Yes, never mind about
        that," she observed. "Yuan Yang's mother also died, as it happens, the other
        day;}
}
{
}

\Verse{17}
{
    \zA{不用说了，我已经猜着了：自然 是王熙凤要求这雏鸾小姐为妻了。” 女先儿笑道：“老祖宗原来听过这回书？” 众
        人都道：“老太太什么没听见过!就是没听见，也猜着了。” 贾母笑道：“这些书 就是一套子，左不过是些佳人才子，最没趣儿。 把人家女儿说的这么坏，还说是‘佳
        人’!编的连影儿也没有了。 开口都是乡绅门第，父亲不是尚书，就是宰相。 一个 小姐，必是爱如珍宝。}
}
{
    \eA{but taking into consideration that both her parents lived in the south, I
        didn't let her return home to observe a period of mourning. But as both
        these girls are now in mourning, why not allow them to live together?
        They'll thus be able to keep each other company. Take a few fruits,
        eatables, and other such things," continuing she bade a matron, "and give
        them to those two girls to eat."}
}
{
}

\Verse{18}
{
    \zA{这小姐必是通文知礼，无所不晓，竟是‘绝代佳人’，只见 了一个清俊男人，不管是亲是友，想起他的终身大事来，父母也忘了，书也忘了，
        鬼不成鬼，贼不成贼，那一点儿像个佳人？ 就是满腹文章，做出这样事来，也算不 得是佳人了。 比如一个男人家，满腹的文章，去做贼，难道那王法看他是个才子就
        不入贼情一案了不成？ 可知那编书的是自己堵自己的嘴。}
}
{
    \eA{"Would she likely wait until now?" Hu Po laughingly interposed. "Why, she
        joined (Hsi Jen) long ago." In the course of this conversation, the various
        inmates partook of some more wine, and watched the theatricals. But we will
        now turn our attention to Pao-yü. He made his way straight into the garden.}
}
{
}

\Verse{19}
{
    \zA{再者：既说是世宦书香大 家子的小姐，又知礼读书，连夫人都知书识礼的，就是告老还家，自然奶妈子丫头
        伏侍小姐的人也不少，怎么这些书上，凡有这样的事，就只小姐和紧跟的一个丫头 知道？ 你们想想，那些人都是管做什么的？ 可是前言不答后语了不是？”
        众人听了，都笑说：“老太太这一说，是谎都批出来了。” 贾母笑道：“有个 原故：编这样书的人，有一等妒人家富贵的，或者有求不遂心，所以编出来遭塌人 家。}
}
{
    \eA{The matrons saw well enough that he was returning to his rooms, but instead
        of following him in, they ensconced themselves near the fire in the tea-room
        situated by the garden-gate, and made the best of the time by drinking and
        playing cards with the girls in charge of the tea. Pao-yü entered the court.
        The lanterns burnt brightly, yet not a human voice was audible. "Have they
        all, forsooth, gone to sleep?"}
}
{
}

\Verse{20}
{
    \zA{再有一等人，他自己看了这些书，看邪了，想着得一个佳人才好，所以编出来 取乐儿。 他何尝知道那世宦读书人家儿的道理!别说那书上那些大家子，如今眼下
        拿着咱们这中等人家说起，也没那样的事。 别叫他诌掉了下巴？ 子罢。 所以我们从 不许说这些书，连丫头们也不懂这些话。
        这几年我老了，他们姐儿们住的远，我偶 然闷了，说几句听听，他们一来，就忙着止住了。” 李薛二人都笑说：“这正是大 家子的规矩。
        连我们家也没有这些杂话叫孩子们听见。”}
}
{
    \eA{She Yüeh ventured. "Let's walk in gently, and give them a fright!"
        Presently, they stepped, on tiptoe, past the mirrored partition-wall. At a
        glance, they discerned Hsi Jen lying on the stove-couch, face to face with
        some other girl. On the opposite side sat two or three old nurses nodding,
        half asleep.}
}
{
}

\Verse{21}
{
    \zA{凤姐儿走上来斟酒，笑道：“罢，罢!酒冷了，老祖宗喝一口润润嗓子再掰谎 罢。 这一回就叫做《掰谎记》，就出在本朝，本地，本年，本月，本日，本时。 老
        祖宗‘一张口难说两家话’，‘花开两朵，各表一枝’，‘是真是谎且不表，再整 观灯看戏的人’。 老祖宗且让这二位亲戚吃杯酒、看两出戏着，再从逐朝话言掰起，
        如何？” 一面说，一面斟酒，一面笑。 未说完，众人俱已笑倒了。}
}
{
    \eA{Pao-yü conjectured that both the girls were plunged in sleep, and was just
        about to enter, when of a sudden some one was heard to heave a sigh and to
        say: "How evident it is that worldly matters are very uncertain! Here you
        lived all alone in here, while your father and mother tarried abroad, and
        roamed year after year from east to west, without any fixed place of abode.}
}
{
}

\Verse{22}
{
    \zA{两个女先儿也笑 个不住，都说：“奶奶好刚口!奶奶要一说书，真连我们吃饭的地方都没了。” 薛 姨妈笑道：“你少兴头些!外头有人，比不得往常。”
        凤姐儿笑道：“外头只有一 位珍大哥哥，我们还是论哥哥妹妹，从小儿一处淘气淘了这么大。}
}
{
    \eA{I ever thought that you wouldn't have been able to be with them at their
        last moments; but, as it happened, (your mother) died in this place this
        year, and you could, after all, stand by her to the end." "Quite so!"
        rejoined Hsi Jen. "Even I little expected to be able to see any of my
        parents' funeral. When I broke the news to our Madame Wang, she also gave me
        forty taels.}
}
{
}

\Verse{23}
{
    \zA{这几年因做了亲， 我如今立了多少规矩了!便不是从小儿兄妹，只论大伯子小婶儿，那二十四孝上‘斑
        衣戏彩’，他们不能来戏彩引老祖宗笑一笑，我这里好容易引的老祖宗笑一笑，多 吃了一点东西，大家喜欢，都该谢我才是，难道反笑我不成？” 贾母笑道：“可是
        这两日我竟没有痛痛的笑一场，倒是亏他才一路说，笑的我这里痛快了些。 我再吃 钟酒。” 吃着酒，又命宝玉：“来敬你姐姐一杯。” 凤姐儿笑道：“不用他敬，我
        讨老祖宗的寿罢。”}
}
{
    \eA{This was really a kind attention on her part. I hadn't nevertheless presumed
        to indulge in any vain hopes." Pao-yü overheard what was said. Hastily
        twisting himself round, he remarked in a low voice, addressing himself to
        She Yüeh and her companions: "Who would have fancied her also in here? But
        were I to enter, she'll bolt away in another tantrum!}
}
{
}

\Verse{24}
{
    \zA{说着便将贾母的杯拿起来，将半杯剩酒吃了，将杯递与丫鬟， 另将温水浸的杯换一个上来。 于是各席上的都撤去，另将温水浸着的代换，斟了新 酒上来，然后归坐。
        女先儿回说：“老祖宗不听这书，或者弹一套曲子听听罢。” 贾母道：“你们 两个对一套《将军令》罢。” 二人听说，忙合弦按调拨弄起来。 贾母因问：“天有
        几更了？” 众婆子忙回：“三更了。” 贾母道：“怪道寒浸浸的起来。” 早有众丫 鬟拿了添换的衣裳送来。}
}
{
    \eA{Better then that we should retrace our steps, and let them quietly have a
        chat together, eh? Hsi Jen was alone, and down in the mouth, so it's a
        fortunate thing that she joined her in such good time." As he spoke, they
        once more walked out of the court with gentle tread. Pao-yü went to the back
        of the rockery, and stopping short, he raised his clothes. She Yüeh and
        Ch'iu Wen stood still, and turned their faces away.}
}
{
}

\Verse{25}
{
    \zA{王夫人起身陪笑说道：“老太太不如挪进暖阁里地炕上， 倒也罢了。 这二位亲戚也不是外人，我们陪着就是了。” 贾母听说，笑道：“既这
        样说，不如大家都挪进去，岂不暖和？” 王夫人道：“恐里头坐不下。” 贾母道： “我有道理：如今也不用这些桌子，只用两三张并起来，大家坐在一处挤着，又亲
        热又暖和。” 众人都道：“这才有趣儿！” 说着，便起了席。 众媳妇忙撤去残席， 里面直顺并了三张大桌，又添换了果馔摆好。}
}
{
    \eA{"Stoop," they smiled, "and then loosen your clothes! Be careful that the
        wind doesn't blow on your stomach!" The two young maids, who followed
        behind, surmised that he was bent upon satisfying a natural want, and they
        hurried ahead to the tea-room to prepare the water. Just, however, as Pao-yü
        was crossing over, two married women came in sight, advancing from the
        opposite direction. "Who's there?" they inquired.}
}
{
}

\Verse{26}
{
    \zA{贾母便说：“都别拘礼，听我分派你 们就坐才好。” 说着，便让薛李正面上坐，自己西向坐了，叫宝琴、黛玉、湘云三
        人皆紧依左右坐下，向宝玉说：“你挨着你太太。” 于是邢夫人王夫人之中夹着宝 玉。 宝钗等姐妹在西边，挨次下去，便是娄氏带着贾蓝、尤氏李纨夹着贾兰,下面
        横头是贾蓉媳妇胡氏。 贾母便说：“珍哥带着你兄弟们去罢，我也就睡了。” 贾珍 等忙答应，又都进来听吩咐。 贾母道：“快去罢，不用进来。
        才坐好了，又都起来。}
}
{
    \eA{"Pao-yü is here," Ch'ing Wen answered. "But mind, if you bawl and shout like
        that, you'll give him a start." The women promptly laughed. "We had no
        idea," they said, "that we were coming, at a great festive time like this,
        to bring trouble upon ourselves! What a lot of hard work must day after day
        fall to your share, young ladies." Speaking the while, they drew near.}
}
{
}

\Verse{27}
{
    \zA{你快歇着罢，明儿还有大事呢。” 贾珍忙答应了，又笑道：“留下蓉儿斟酒才是。” 贾母笑道：“正是忘了他。” 贾珍应了一个“是”，便转身带领贾琏等出来。
        二人 自是欢喜，便命人将贾琮贾璜各自送回家去，便约了贾琏去追欢买笑，不在话下。 这里贾母笑道：“我正想着，虽然这些人取乐，必得重孙一对双全的在席上才
        好。 蓉儿这可全了。 蓉儿!和你媳妇坐在一处，倒也团圆了。” 因有家人媳妇呈上 戏单，贾母笑道：“我们娘儿们正说得兴头，又要吵起来。}
}
{
    \eA{She Yüeh and her friends then asked them what they were holding in their
        hands. "We're taking over," they replied, "some things to the two girls:
        Miss Chin and Miss Hua." "They're still singing the 'Eight Worthies'
        outside," She Yüeh went on to observe laughingly, "and how is it you're
        running again to Miss Chin's and Miss Hua's before the 'Trouble-first
        moon-box' has been gone through?"}
}
{
}

\Verse{28}
{
    \zA{况且那孩子们熬夜，怪 冷的。 也罢，且叫他们歇歇，把咱们的女孩子们叫起来，就在这台上唱两出罢，也 给他们瞧瞧。”
        媳妇子们听了，答应出来，忙的一面着人往大观园去传人，一面二 门口去传小厮们伺候。 小厮们忙至戏房，将班中所有大人一概带出，只留下小孩子 们。
        一时，梨香院的教习带了文官等十二人从游廊角门出来，婆子们抱着几个软 包，因不及抬箱，料着贾母爱听的三五出戏的彩衣包了来。 婆子们带了文官等进去，
        见过，只垂手站着。}
}
{
    \eA{"Take the lid off," Pao-yü cried, "and let me see what there's inside."
        Ch'in Wen and She Yüeh at once approached and uncovered the boxes.}
}
{
}

\Verse{29}
{
    \zA{贾母笑道：“大正月里，你师父也不放你们出来逛逛？ 你们如 今唱什么？ 才刚八出《八义》，闹的我头疼，咱们清淡些好。 你瞧瞧，薛姨太太，
        这李亲家太太，都是有戏的人家，不知听过多少好戏的； 这些姑娘们都比咱们家的 姑娘见过好戏，听过好曲子。 如今这小戏子又是那有名玩戏的人家的班子，虽是小
        孩子，却比大班子还强。 咱们好歹别落了褒贬!少不得弄个新样儿的：叫芳官唱一 出《寻梦》，只用箫和笙笛，馀者一概不用。” 文官笑道：“老祖宗说的是。}
}
{
    \eA{The two women promptly stooped, which enabled Pao-yü to see that the
        contents of the two boxes consisted alike of some of the finest fruits and
        tea-cakes, which had figured at the banquet, and, nodding his head, he
        walked off, while She Yüeh and her friend speedily threw the lids down
        anyhow, and followed in his track. "Those two dames are pleasant enough,"
        Pao-yü smiled, "and they know how to speak decently;}
}
{
}

\Verse{30}
{
    \zA{我们 的戏，自然不能入姨太太和亲家太太姑娘们的眼； 不过听我们一个发脱口齿，再听 个喉咙罢了。” 贾母笑道：“正是这话了。”
        李婶娘薛姨妈喜的笑道：“好个灵透 孩子，你也跟着老太太打趣我们。” 贾母笑道：“我们这原是随便的玩意儿，又不 出去做买卖，所以竟不大合时。”
        说着，又叫葵官：“唱一出《惠明下书》，也不 用抹脸。 只用这两出，叫他们二位太太听个助意儿罢了。 若省了一点儿力，我可不 依。”}
}
{
    \eA{but it's they who get quite worn out every day, and they contrariwise say
        that you've got ample to do daily. Now, doesn't this amount to bragging and
        boasting?" "Those two women," She Yüeh chimed in, "are not bad. But such of
        them as don't know what good manners mean are ignorant to a degree of all
        propriety." "You, who know what's what," Pao-yü added, "should make
        allowances for that kind of rustic people.}
}
{
}

\Verse{31}
{
    \zA{文官等听了出来，忙去扮演上台，先是《寻梦》，次是《下书》。 众人鸦雀 无闻。 薛姨妈笑道：“实在戏也看过几百班，从没见过只用箫管的。” 贾母道：“先
        有，只是像方才《西楼》《楚江情》一只，多有小生吹箫合的。 这合大套的实在少。 这也在人讲究罢了，这算什么出奇。” 又指着湘云道：“我像他这么大的时候儿，
        他爷爷有一班小戏，偏有一个弹琴的，凑了《西厢记》的《听琴》，《玉簪记》的 《琴挑》，《续琵琶》的《胡茄十八拍》，竟成了真的了。}
}
{
    \eA{You should pity them; that's all." Speaking, he made his exit out of the
        garden gate. The matrons had, though engaged in drinking and gambling, kept
        incessantly stepping out of doors to furtively keep an eye on his movements,
        so that the moment they perceived Pao-yü appear, they followed him in a
        body. On their arrival in the covered passage of the reception-hall, they
        espied two young waiting-maids;}
}
{
}

\Verse{32}
{
    \zA{比这个更如何？” 众人 都道：“那更难得了。” 贾母于是叫过媳妇们来，吩咐文官等叫他们吹弹一套《灯 月圆》。 媳妇们领命而去。
        当下贾蓉夫妻二人捧酒一巡。 凤姐儿因贾母十分高兴，便笑道：“趁着女先儿 们在这里，不如咱们传梅，行一套‘春喜上眉梢’的令，如何？” 贾母笑道：“这
        是个好令啊!正对时景儿。” 忙命人取了黑漆铜钉花腔令鼓来，给女先儿击着。 席 上取了一枝红梅，贾母笑道：“到了谁手里住了鼓，吃一杯，也要说些什么才好。”}
}
{
    \eA{the one with a small basin in her hand; the other with a towel thrown over
        her arm. They also held a bowl and small kettle, and had been waiting in
        that passage for ever so long. Ch'iu Wen was the first to hastily stretch
        out her hand and test the water. "The older you grow," she cried, "the
        denser you get! How could one ever use this icy-cold water?" "Miss, look at
        the weather!" the young maid replied.}
}
{
}

\Verse{33}
{
    \zA{凤姐儿笑道：“依我说，谁像老祖宗要什么有什么呢？ 我们这不会的不没意思吗？ 怎 么能雅俗共赏才好。 不如谁住了，谁说个笑话儿罢。”
        众人听了，都知道他素日善 说笑话儿，肚内有无限的新鲜趣令； 今见如此说，不但在席的诸人喜欢，连地下伏 侍的老小人等无不欢喜。
        那小丫头子们都忙去找姐姐叫妹妹的，告诉他们：“快来 听，二奶奶又说笑话儿了。” 众丫头子们便挤了一屋子。
        于是戏完乐罢，贾母将些汤细点果给文官等吃去，便命响鼓。}
}
{
    \eA{"I was afraid the water would get cold. It was really scalding; is it cold
        now?" While she made this rejoinder, an old matron was, by a strange
        coincidence, seen coming along, carrying a jug of hot water. "Dear dame,"
        shouted the young maid, "come over and pour some for me in here!" "My dear
        girl," the matron responded, "this is for our old mistress to brew tea with.
        I'll tell you what;}
}
{
}

\Verse{34}
{
    \zA{那女先儿们都是 惯熟的，或紧或慢，或如残漏之滴，或如迸豆之急，或如惊马之驰，或如疾电之光， 忽然暗其鼓声，那梅方递至贾母手中，鼓声恰住，大家哈哈大笑。
        贾蓉忙上来斟了 一杯，众人都笑道：“自然老太太先喜了，我们才托赖些喜。” 贾母笑道：“这酒 也罢了，只是这笑话儿倒有些难说。”
        众人都说：“老太太的比凤姑娘说的还好， 赏一个，我们也笑一笑。” 贾母笑道：“并没有新鲜招笑儿的，少不得老脸皮厚的 说一个罢。”}
}
{
    \eA{you'd better go and fetch some yourself. Are you perchance afraid lest your
        feet might grow bigger by walking?" "I don't care whose it is," Ch'iu Wen
        put in. "If you don't give me any, I shall certainly empty our old lady's
        teapot and wash my hands." The old matron turned her head; and, catching
        sight of Ch'iu Wen, she there and then raised the jug and poured some of the
        water. "That will do!" exclaimed Ch'iu Wen.}
}
{
}

\Verse{35}
{
    \zA{因说道： “一家子养了十个儿子，娶了十房媳妇儿。 惟有第十房媳妇儿聪明伶俐、心巧 嘴乖，公婆最疼，成日家说那九个不孝顺。
        这九个媳妇儿委屈，便商议说：‘咱们 九个心里孝顺，只是不像那小蹄子儿嘴巧，所以公公婆婆只说他好。 这委屈向谁诉 去？’
        有主意的说道：‘咱们明儿到阎王庙去烧香，和阎王爷说去，问他一问：叫 我们托生为人，怎么单单给那小蹄子儿一张乖嘴，我们都入了夯嘴里头？’ 那八个
        听了，都喜欢说：‘这个主意不错。’}
}
{
    \eA{"With all your years, don't you yet know what's what? Who isn't aware that
        it's for our old mistress? But would one presume to ask for what shouldn't
        be asked for?" "My eyes are so dim," the matron rejoined with a smile, "that
        I didn't recognise this young lady." When Pao-yü had washed his hands, the
        young maid took the small jug and filled the bowl; and, as she held it in
        her hand, Pao-yü rinsed his mouth.}
}
{
}

\Verse{36}
{
    \zA{第二日，便都往阎王庙里来烧香。 九个都在 供桌底下睡着了。 九个魂专等阎王驾到。 左等不来，右等也不到。 正着急，只见孙
        行者驾着斤斗云来了，看见九个魂，便要拿金箍棒打来。 吓得九个魂忙跪下央求。 孙行者问起原故来，九个人忙细细的告诉了他。 孙行者听了，把脚一跺，叹了一口
        气道：‘这原故幸亏遇见我!等着阎王来了，他也不得知道。’ 九个人听了，就求 说：‘大圣发个慈悲，我们就好了。’}
}
{
    \eA{But Ch'iu Wen and She Yüeh availed themselves likewise of the warm water to
        have a wash; after which, they followed Pao-yü in. Pao-yü at once asked for
        a kettle of warm wine, and, starting from sister-in-law Li, he began to
        replenish their cups. (Sister-in-law Li and his aunt Hsüeh) pressed him,
        however, with smiling faces, to take a seat; but his grandmother Chia
        remonstrated.}
}
{
}

\Verse{37}
{
    \zA{孙行者笑道：‘却也不难：那日你们妯娌十 个托生时，可巧我到阎王那里去，因为撒了一泡尿在地下，你那个小婶儿便吃了。
        你们如今要伶俐嘴乖，有的是尿，便撒泡你们吃就是了。’” 说毕，大家都笑起来。 凤姐儿笑道：“好的呀!幸而我们都是夯嘴夯腮的，不 然，也就吃了猴儿尿了！”
        尤氏娄氏都笑向李纨道：“咱们这里头谁是吃过猴儿尿 的，别装没事人儿！” 薛姨妈笑道：“笑话儿在对景就发笑。” 说着，又击起鼓来。}
}
{
    \eA{"He's only a youngster," she said, "so let him pour the wine! We must all
        drain this cup!" With these words, she quaffed her own cup, leaving no
        heel-taps. Mesdames Hsing and Wang also lost no time in emptying theirs; so
        Mrs. Hsüeh and 'sister-in-law' Li had no alternative but to drain their
        share. "Fill the cups too of your female cousins, senior or junior," dowager
        lady Chia went on to tell Pao-yü.}
}
{
}

\Verse{38}
{
    \zA{小丫头子们只要听凤姐儿的笑话，便悄悄的和女先儿说明， 以咳嗽为记。 须臾传至两遍，刚到了凤姐儿手里，小丫头子们故意咳嗽，女先儿便 住了。
        众人齐笑道：“这可拿住他了!快吃了酒，说一个好的罢，别太逗人笑的肠 子疼！” 凤姐儿想一想，笑道：“一家子也是过正月节，合家赏灯吃酒，真真的热闹非 常。}
}
{
    \eA{"And you mayn't pour the wine anyhow. Each of you must swallow every drop of
        your drinks." Pao-yü upon hearing her wishes, set to work, while signifying
        his assent, to replenish the cups of the several young ladies in their
        proper gradation. But when he got to Tai-yü, she raised the cup, for she
        would not drink any wine herself, and applied it to Pao-yü's lips. Pao-yü
        drained the contents with one breath;}
}
{
}

\Verse{39}
{
    \zA{祖婆婆、太婆婆、媳妇、孙子媳妇、重孙子媳妇、亲孙子媳妇、侄孙子、重孙 子、灰孙子、滴里搭拉的孙子、孙女儿、外孙女儿、姨表孙女儿、姑表孙女儿……
        嗳哟哟!真好热闹！” 众人听他说着，已经笑了，都说：“听这数贫嘴的!又不知要 编派那一个呢！” 尤氏笑道：“你要招我，我可撕你的嘴！”
        凤姐儿起身拍手笑道： “人家这里费力，你们紧着混，我就不说了。” 贾母笑道：“你说你的，底下怎么 样？”}
}
{
    \eA{upon which Tai-yü gave him a smile, and said to him: "I am much obliged to
        you." Pao-yü next poured a cup for her. But lady Feng immediately laughed
        and expostulated. "Pao-yü!" she cried, "you mustn't take any cold wine.
        Mind, your hand will tremble, and you won't be able to-morrow to write your
        characters or to draw the bow." "I'm not having any cold wine," Pao-yü
        replied.}
}
{
}

\Verse{40}
{
    \zA{凤姐儿想了一想，笑道：“底下就团团的坐了一屋子，吃了一夜酒，就散了。” 众人见他正言厉色的说了，也都再无有别话，怔怔的还等往下说，只觉他冰冷
        无味的就住了。 湘云看了他半日。 凤姐儿笑道：“再说一个过正月节的：几个人拿 着房子大的炮仗往城外放去，引了上万的人跟着瞧去。 有一个性急的人等不得，就
        偷着拿香点着了。 只见‘噗哧’的一声，众人哄然一笑，都散了。 这抬炮仗的人抱 怨卖炮仗的捍的不结实，没等放就散了。”}
}
{
    \eA{"I know you're not," lady Feng smiled, "but I simply warn you." After this,
        Pao-yü finished helping the rest of the inmates inside, with the exception
        of Chia Jung's wife, for whom he bade a maid fill a cup. Then emerging again
        into the covered passage, he replenished the cups of Chia Chen and his
        companions; after which, he tarried with them for a while, and at last
        walked in and resumed his former seat.}
}
{
}

\Verse{41}
{
    \zA{湘云道：“难道本人没听见？” 凤姐儿 道：“本人原是个聋子。” 众人听说，想了一回，不觉失声都大笑起来。 又想着先
        前那个没完的，问他道：“先那一个到底怎么样？ 也该说完了。” 凤姐儿将桌子一 拍，道：“好罗唆!到了第二日是十六日，年也完了，节也完了，我看人忙着收东
        西还闹不清，那里还知道底下的事了？” 众人听说，复又笑起。 凤姐儿笑道：“外头已经四更多了，依我说：老祖宗也乏了，咱们也该‘聋子 放炮仗――散了’罢？”}
}
{
    \eA{Presently, the soup was brought, and soon after that the 'feast of lanterns'
        cakes were handed round. Dowager lady Chia gave orders that the play should
        be interrupted for a time. "Those young people," (she said) "are be to
        pitied! Let them too have some hot soup and warm viands. They then can go on
        again.}
}
{
}

\Verse{42}
{
    \zA{尤氏等用绢子握着嘴，笑的前仰后合，指他说道：“这个 东西真会数贫嘴！” 贾母笑道：“真真这凤丫头，越发炼贫了！” 一面说，一面吩
        咐道：“他提起炮仗来，咱们也把烟火放了，解解酒。” 贾蓉听了，忙出去带着小 厮们就在院子内安下屏架，将烟火设吊齐备。 这烟火俱系各处进贡之物，虽不甚大，
        却极精致，各色故事俱全，夹着各色的花炮。 黛玉禀气虚弱，不禁劈拍之声，贾母 便搂他在怀内。 薛姨妈便搂湘云，湘云笑道：“我不怕。”}
}
{
    \eA{Take of every kind of fruit," she continued, "'feast of lanterns' cakes, and
        other such dainties and give them a few." The play was shortly stopped. The
        matrons ushered in a couple of blind singing-girls, who often came to the
        house, and put two benches, on the opposite side, for them. Old lady Chia
        desired them to take a seat, and banjos and guitars were then handed to
        them. "What stories would you like to hear?"}
}
{
}

\Verse{43}
{
    \zA{宝钗笑道：“他专爱自 己放大炮仗，还怕这个呢！” 王夫人便将宝玉搂入怀内。 凤姐笑道：“我们是没人 疼的！” 尤氏笑道：“有我呢，我搂着你。
        ――你这会子又撒娇儿了，听见放炮仗， 就像‘吃了蜜蜂儿屎’的，今儿又轻狂了。” 凤姐儿笑道：“等散了，咱们园子里 放去，我比小厮们还放的好呢。”
        说话之间，外面一色色的放了又放。 又有许多“满 天星”“九龙入云”“平地一声雷”“飞天十响”之类的零星小炮仗。}
}
{
    \eA{old lady Chia inquired of 'sister-in-law' Li and Mrs. Hsüeh. "We don't care
        what they are;" both of them rejoined with one voice. "Any will do!" "Have
        you of late added any new stories to your stock?" old lady Chia asked.
        "We've got a new story," the two girls explained. "It's about an old affair
        of the time of the Five Dynasties, which trod down the T'ang dynasty."
        "What's its title?" old lady Chia inquired.}
}
{
}

\Verse{44}
{
    \zA{放罢，然后 又命小戏子打了一回“莲花落”，撒得满台的钱，那些孩子们满台的抢钱取乐。 上汤时，贾母说：“夜长，不觉得有些饿了。”
        凤姐忙回说：“有预备的鸭子 肉粥。” 贾母道：“我吃些清淡的罢。” 凤姐儿忙道：“也有枣儿熬的粳米粥，预 备太太们吃斋的。”
        贾母道：“倒是这个还罢了。” 说着，已经撤去残席，内外另 设各种精致小菜。 大家随意吃了些，用过漱口茶，方散。
        十七日一早，又过宁府行礼，伺候掩了祠门，收过影像，方回来。}
}
{
    \eA{"It's called: 'A Feng seeks a Luan in marriage': (the male phoenix asks the
        female phoenix in marriage)," one of the girls answered. "The title is all
        very well," dowager lady Chia proceeded, "but why I wonder was it ever given
        to it. First tell us its general purport, and if it's interesting, you can
        continue." "This story," the girl explained, "treats of the time when the
        T'ang dynasty was extinguished.}
}
{
}

\Verse{45}
{
    \zA{此日便是薛 姨妈家请吃年酒。 贾母连日觉得身上乏了，坐了半日，回来了。 自十八日以后，亲
        友来请或来赴席的，贾母一概不会，有邢夫人、王夫人、凤姐三人料理。 连宝玉只 除王子腾家去了，馀者亦皆不去，只说是贾母留下解闷。
        当下元宵已过，凤姐忽然小产了，合家惊慌。 要知端底，下回分解。 (本章完)}
}
{
    \eA{There lived then one of the gentry, who had originally been a denizen of
        Chin Ling. His name was Wang Chun. He had been minister under two reigns. He
        had, about this time, pleaded old age and returned to his home. He had about
        his knees only one son, called Wang Hsi-feng." When the company heard so
        far, they began to laugh. "Now isn't this a duplicate of our girl Feng's
        name?" old lady Chia laughingly exclaimed.}
}
{
}

\Verse{46}
{
    \zA{}
}
{
    \eA{A married woman hurried up and pushed (the girl). "That's the name of your
        lady Secunda," she said, "so don't use it quite so heedlessly!" "Go on with
        your story!" dowager lady Chia shouted. The girl speedily stood up, smiling
        the while. "We do deserve death!" she observed. "We weren't aware that it
        was our lady's worthy name." "Why should you be in such fear and trembling?"
        lady Feng laughed. "Go on! There are many duplicate names and duplicate
        surnames." The girl then proceeded with her story. "In a certain year," she
        resumed, "his honour old Mr. Wang saw his son Mr. Wang off for the capital
        to be in time for the examinations. One day, he was overtaken by a heavy
        shower of rain and he betook himself into a village for shelter.}
}
{
}

\Verse{47}
{
    \zA{}
}
{
    \eA{Who'd have thought it, there lived in this village, one of the gentry, of
        the name of Li, who had been an old friend of his honour old Mr. Wang, and
        he kept Mr. Wang junior to put up in his library. This Mr. Li had no son,
        but only a daughter. This young daughter's worthy name was Ch'u Luan. She
        could perform on the lute; she could play chess; and she had a knowledge of
        books and of painting. There was nothing that she did not understand." Old
        lady Chia eagerly chimed in. "It's no wonder," she said, "that the story has
        been called: 'A Feng seeks a Luan in marriage,' '(a male phoenix seeks a
        female phoenix in marriage).' But you needn't proceed. I've already guessed
        the denouement. There's no doubt that Wang Hsi-feng asks for the hand of
        this Miss Ch'u Luan."}
}
{
}

\Verse{48}
{
    \zA{}
}
{
    \eA{"Your venerable ladyship must really have heard the story before," the
        singing-girl smiled. "What hasn't our worthy senior heard?" they all
        exclaimed. "But she's quick enough in guessing even unheard of things." "All
        these stories run invariably in one line," old lady Chia laughingly
        rejoined. "They're all about pretty girls and scholars. There's no fun in
        them. They abuse people's daughters in every possible way, and then they
        still term them nice pretty girls. They're so concocted that there's not
        even a semblance of truth in them. From the very first, they canvass the
        families of the gentry. If the paterfamilias isn't a president of a board;
        then he's made a minister. The heroine is bound to be as lovable as a gem.}
}
{
}

\Verse{49}
{
    \zA{}
}
{
    \eA{This young lady is sure to understand all about letters, and propriety. She
        knows every thing and is, in a word, a peerless beauty. At the sight of a
        handsome young man, she pays no heed as to whether he be relation or friend,
        but begins to entertain thoughts of the primary affair of her life, and
        forgets her parents and sets her books on one side. She behaves as neither
        devil nor thief would: so in what respect does she resemble a nice pretty
        girl? Were even her brain full of learning, she couldn't be accounted a nice
        pretty girl, after behaving in this manner! Just like a young fellow, whose
        mind is well stored with book-lore, and who goes and plays the robber!}
}
{
}

\Verse{50}
{
    \zA{}
}
{
    \eA{Now is it likely that the imperial laws would look upon him as a man of
        parts, and that they wouldn't bring against him some charge of robbery? From
        this it's evident that those, who fabricate these stories, contradict
        themselves. Besides, they may, it's true, say that the heroines belong to
        great families of official and literary status, that they're conversant with
        propriety and learning and that their honourable mothers too understand
        books and good manners, but great households like theirs must, in spite of
        the parents having pleaded old age and returned to their natives places,
        contain a great number of inmates;}
}
{
}

\Verse{51}
{
    \zA{}
}
{
    \eA{and the nurses, maids and attendants on these young ladies must also be
        many; and how is it then that, whenever these stories make reference to such
        matters, one only hears of young ladies with but a single close attendant?
        What can, think for yourselves, all the other people be up to? Indeed, what
        is said before doesn't accord with what comes afterwards. Isn't it so, eh?"
        The party listened to her with much glee. "These criticisms of yours,
        venerable ancestor," they said, "have laid bare every single discrepancy."
        "They have however their reasons," old lady Chia smilingly resumed.}
}
{
}

\Verse{52}
{
    \zA{}
}
{
    \eA{"Among the writers of these stories, there are some, who begrudge people's
        wealth and honours, or possibly those, who having solicited a favour (of the
        wealthy and honorable), and not obtained the object, upon which their wishes
        were set, have fabricated lies in order to disparage people. There is
        moreover a certain class of persons, who become so corrupted by the perusal
        of such tales that they are not satisfied until they themselves pounce upon
        some nice pretty girl. Hence is it that, for fun's sake, they devise all
        these yarns. But how could such as they ever know the principle which
        prevails in official and literary families?}
}
{
}

\Verse{53}
{
    \zA{}
}
{
    \eA{Not to speak of the various official and literary families spoken about in
        these anecdotes, take now our own immediate case as an instance. We're only
        such a middle class household, and yet we've got none of those occurrences;
        so don't let her go on spinning these endless yarns. We must on no account
        have any of these stories told us! Why, even the maids themselves don't
        understand any of this sort of language. I've been getting so old the last
        few years, that I felt unawares quite melancholy whenever the girls went to
        live far off, so my wont has been to have a few passages recounted to me;
        but as soon as they got back, I at once put a stop to these things."
        'Sister-in-law' Li and Mrs. Hsüeh both laughed. "This is just the rule,"
        they said, "which should exist in great families.}
}
{
}

\Verse{54}
{
    \zA{}
}
{
    \eA{Not even in our homes is any of this confused talk allowed to reach the ears
        of the young people." Lady Feng came forward and poured some wine. "Enough,
        that will do!" she laughed. "The wine has got quite cold. My dear ancestor,
        do take a sip and moisten your throat with, before you begin again to dilate
        on falsehoods. What we've been having now can well be termed 'Record of a
        discussion on falsehoods.' It has had its origin in this reign, in this
        place, in this year, in this moon, on this day and at this very season. But,
        venerable senior, you've only got one mouth, so you couldn't very well
        simultaneously speak of two families. 'When two flowers open together,' the
        proverb says, 'one person can only speak of one.' But whether the stories be
        true or fictitious, don't let us say anything more about them.}
}
{
}

\Verse{55}
{
    \zA{}
}
{
    \eA{Let's have the footlights put in order, and look at the players. Dear
        senior, do let these two relatives have a glass of wine and see a couple of
        plays; and you can then start arguing about one dynasty after another. Eh,
        what do you say?" Saying this, she poured the wine, laughing the while. But
        she had scarcely done speaking before the whole company were convulsed with
        laughter. The two singing girls were themselves unable to keep their
        countenance. "Lady Secunda," they both exclaimed, "what a sharp tongue you
        have! Were your ladyship to take to story-telling, we really would have
        nowhere to earn our rice." "Don't be in such overflowing spirits," Mrs.
        Hsüeh laughed. "There are people outside; this isn't like any ordinary
        occasion."}
}
{
}

\Verse{56}
{
    \zA{}
}
{
    \eA{"There's only my senior brother-in-law Chen outside," lady Feng smiled. "And
        we've been like brother and sister from our youth up. We've romped and been
        up to every mischief to this age together. But all on account of my
        marriage, I've had of late years to stand on ever so many ceremonies. Why
        besides being like brother and sister from the time we were small kids, he's
        anyhow my senior brother-in-law, and I his junior sister-in-law. (One among)
        those twenty four dutiful sons, travestied himself in theatrical costume (to
        amuse his parents), but those fellows haven't sufficient spirit to come in
        some stage togs and try and make you have a laugh, dear ancestor.}
}
{
}

\Verse{57}
{
    \zA{}
}
{
    \eA{I've however succeeded, after ever so much exertion, in so diverting you as
        to induce you to eat a little more than you would, and in putting everybody
        in good humour; and I should be thanked by one and all of you; it's only
        right that I should. But can it be that you will, on the contrary, poke fun
        at me?" "I've truly not had a hearty laugh the last few days," old lady Chia
        smiled, "but thanks to the funny things she recounted just now, I've managed
        to get in somewhat better spirits in here. So I'll have another cup of
        wine." Then having drunk her wine, "Pao-yü," she went on to say, "come and
        present a cup to your sister-in-law!" Lady Feng gave a smile. "There's no
        use for him to give me any wine," she ventured.}
}
{
}

\Verse{58}
{
    \zA{}
}
{
    \eA{"(I'll drink out of your cup,) so as to bring upon myself your longevity,
        venerable ancestor." While uttering this response, she raised dowager lady
        Chia's cup to her lips, and drained the remaining half of the contents;
        after which, she handed the cup to a waiting-maid, who took one from those
        which had been rinsed with tepid water, and brought it to her. But in due
        course, the cups from the various tables were cleared, and clean ones,
        washed in warm water, were substituted; and when fresh wine had been served
        round, (lady Feng and the maid) resumed their seats. "Venerable lady," a
        singing-girl put in, "you don't like the stories we tell; but may we thrum a
        song for you?"}
}
{
}

\Verse{59}
{
    \zA{}
}
{
    \eA{"You two," remarked old lady Chia, "had better play a duet of the 'Chiang
        Chün ling' song: 'the general's command.'" Hearing her wishes, the two girls
        promptly tuned their cords, to suit the pitch of the song, and struck up on
        their guitars. "What watch of the night is it?" old lady Chia at this point
        inquired. "It's the third watch," the matrons replied with alacrity. "No
        wonder it has got so chilly and damp!" old lady Chia added. Extra clothes
        were accordingly soon fetched by the servants and maids. Madame Wang
        speedily rose to her feet and forced a smile. "Venerable senior," she said,
        "wouldn't it be prudent for you to move on to the stove couch in the winter
        apartments? It would be as well. These two relatives are no strangers. And
        if we entertain them, it will be all right."}
}
{
}

\Verse{60}
{
    \zA{}
}
{
    \eA{"Well, in that case," dowager lady Chia smilingly rejoined, "why shouldn't
        the whole company adjourn inside? Wouldn't it be warmer for us all?" "I'm
        afraid there isn't enough sitting room for every one of us," Madame Wang
        explained. "I've got a plan," old lady Chia added. "We can now dispense with
        these tables. All we need are two or three, placed side by side; we can then
        sit in a group, and by bundling together it will be both sociable as well as
        warm." "Yes, this will be nice!" one and all cried. Assenting, they
        forthwith rose from table. The married women hastened to remove the
        debandade of the banquet.}
}
{
}

\Verse{61}
{
    \zA{}
}
{
    \eA{Then placing three large tables lengthways side by side in the inner rooms,
        they went on to properly arrange the fruits and viands, some of which had
        been replenished, others changed. "You must none of you stand on any
        ceremonies!" dowager lady Chia observed. "If you just listen while I allot
        you your places, and sit down accordingly, it will be all right!"
        Continuing, she motioned to Mrs. Hsüeh and 'sister-in-law' Li to take the
        upper seats on the side of honour, and, making herself comfortable on the
        west, she bade the three cousins Pao-ch'in, Tai-yü and Hsian-yün sit close
        to her on the left and on the right. "Pao-yü," she proceeded "you must go
        next to your mother."}
}
{
}

\Verse{62}
{
    \zA{}
}
{
    \eA{So presently she put Pao-yü, and Pao-ch'ai and the rest of the young ladies
        between Mesdames Hsing and Wang. On the west, she placed, in proper
        gradation, dame Lou, along with Chia Lan, and Mrs. Yu and Li Wan, with Chia
        Lan, (number two,) between them. While she assigned a chair to Chia Jung's
        wife among the lower seats, put crosswise. "Brother Chen," old lady Chia
        cried, "take your cousins and be off! I'm also going to sleep in a little
        time." Chia Chen and his associates speedily expressed their obedience, and
        made, in a body, their appearance inside again to listen to any injunctions
        she might have to give them. "Bundle yourself away at once!" shouted dowager
        lady Chia. "You needn't come in. We've just sat down, and you'll make us get
        up again. Go and rest; be quick!}
}
{
}

\Verse{63}
{
    \zA{}
}
{
    \eA{To-morrow, there are to be some more grand doings!" Chia Chen assented with
        alacrity. "But Jung Erh should remain to replenish the cups," he smiled;
        "it's only fair that he should." "Quite so!" answered old lady Chia
        laughingly. "I forgot all about him." "Yes!" acquiesced Chia Chen. Then
        twisting himself round, he led Chia Lien and his companions out of the
        apartment. (Chia Chen and Chia Lien) were, of course, both pleased at being
        able to get away. So bidding the servants see Chia Tsung and Chia Huang to
        their respective homes, (Chia Chen) arranged with Chia Lien to go in pursuit
        of pleasure and in quest of fun. But we will now leave them to their own
        devices without another word.}
}
{
}

\Verse{64}
{
    \zA{}
}
{
    \eA{"I was just thinking," meanwhile dowager lady Chia laughed, "that it would
        be well, although you people are numerous enough to enjoy yourselves, to
        have a couple of great-grandchildren present at this banquet, so Jung Erh
        now makes the full complement. But Jung Erh sit near your wife, for she and
        you will then make the pair complete." The wife of a domestic thereupon
        presented a play-bill. "We, ladies," old lady Chia demurred, "are now
        chatting in high glee, and are about to start a romp. Those young folks
        have, also, been sitting up so far into the night that they must be quite
        cold, so let the plays alone. Tell them then to have a rest. Yet call our
        own girls to come and sing a couple of plays on this stage. They too will
        thus have a chance of watching us a bit."}
}
{
}

\Verse{65}
{
    \zA{}
}
{
    \eA{After lending an ear to her, the married women assented and quitted the
        room. And immediately finding some servant to go to the garden of Broad
        Vista and summon the girls, they betook themselves, at the same time, as far
        as the second gate and called a few pages to wait on them. The pages went
        with hurried step to the rooms reserved for the players, and taking with
        them the various grown-up members of the company, they only left the more
        youthful behind. Then fetching, in a little time, Wen Kuan and a few other
        girls, twelve in all, from among the novices in the Pear Fragrance court,
        they egressed by the corner gate leading out of the covered passage.}
}
{
}

\Verse{66}
{
    \zA{}
}
{
    \eA{The matrons took soft bundles in their arms, as their strength was not equal
        to carrying boxes. And under the conviction that their old mistress would
        prefer plays of three or five acts, they had put together the necessary
        theatrical costumes. After Wen Kuan and the rest of the girls had been
        introduced into the room by the matrons, they paid their obeisance, and,
        dropping their arms against their sides, they stood reverentially. "In this
        propitious first moon," old lady Chia smiled, "won't your teacher let you
        come out for a stroll? What are you singing now?}
}
{
}

\Verse{67}
{
    \zA{}
}
{
    \eA{The eight acts of the 'Eight worthies' recently sung here were so noisy,
        that they made my head ache; so you'd better let us have something more
        quiet. You must however bear in mind that Mrs. Hsüeh and Mrs. Li are both
        people, who give theatricals, and have heard I don't know how many fine
        plays. The young ladies here have seen better plays than our own girls; and
        they have heard more beautiful songs than they. These actresses, you see
        here now, formed once, despite their youth, part of a company belonging to
        renowned families, fond of plays; and though mere children, they excel any
        troupe composed of grown-up persons. So whatever we do, don't let us say
        anything disparaging about them. But we must now have something new. Tell
        Fang Kuan to sing us the 'Hsün Meng' ballad; and let only flutes and Pandean
        pipes be used.}
}
{
}

\Verse{68}
{
    \zA{}
}
{
    \eA{The other instruments can be dispensed with." "Your venerable ladyship is
        quite right," Wen Kuan smiled. "Our acting couldn't, certainly, suit the
        taste of such people as Mrs. Hsüeh, Mrs. Li and the young ladies.
        Nevertheless, let them merely heed our enunciation, and listen to our
        voices; that's all." "Well said!" dowager lady Chia laughed. 'Sister-in-law'
        Li and Mrs. Hsüeh were filled with delight. "What a sharp girl!" they
        remarked smilingly. "But do you also try to imitate our old lady by pulling
        our leg?" "They're intended to afford us some ready-at-hand recreation," old
        lady Chia smiled. "Besides, they don't go out to earn money. That's how it
        is they are not so much up to the times." At the close of this remark, she
        also desired K'uei Kuan to sing the play: 'Hui Ming sends a letter.'}
}
{
}

\Verse{69}
{
    \zA{}
}
{
    \eA{"You needn't," she added, "make your face up. Just sing this couple of plays
        so as to merely let both those ladies hear a kind of parody of them. But if
        you spare yourselves the least exertion, I shall be unhappy." When they
        heard this, Wen Kuan and her companions left the apartment and promptly
        apparelled themselves and mounted the stage. First in order, was sung the
        'Hsün Meng;' next, '(Hui Ming) sends a letter;' during which, everybody
        observed such perfect silence that not so much as the caw of a crow fell on
        the ear. "I've verily seen several hundreds of companies," Mrs. Hsüeh
        smiled, "but never have I come across any that confined themselves to
        flutes." "There are some," dowager lady Chia answered.}
}
{
}

\Verse{70}
{
    \zA{}
}
{
    \eA{"In fact, in that play acted just now called: 'Love in the western tower at
        Ch'u Ch'iang,' there's a good deal sung by young actors in unison with the
        flutes. But lengthy unison pieces of this description are indeed few. This
        too, however, is purely a matter of taste; there's nothing out of the way
        about it. When I was of her age," resuming, she pointed at Hsiang-yün, "her
        grandfather kept a troupe of young actresses.}
}
{
}

\Verse{71}
{
    \zA{}
}
{
    \eA{There was among them one, who played the lute so efficiently that she
        performed the part when the lute is heard in the 'Hsi Hsiang Chi,' the piece
        on the lute in the 'Yü Ts'an Chi,' and that in the supplementary 'P'i Pa
        Chi,' on the Mongol flageolet with the eighteen notes, in every way as if
        she had been placed in the real circumstances herself. Yea, far better than
        this!" "This is still rarer a thing!" the inmates exclaimed. Old lady Chia
        then shortly called the married women, and bade them tell Wen Kúan and the
        other girls to use both wind and string instruments and render the piece;
        'At the feast of lanterns, the moon is round.' The women servants received
        her orders and went to execute them. Chia Jung and his wife meanwhile passed
        the wine round.}
}
{
}

\Verse{72}
{
    \zA{}
}
{
    \eA{When lady Feng saw dowager lady Chia in most exuberant spirits, she smiled.
        "Won't it be nice," she said, "to avail ourselves of the presence of the
        singing girls to pass plum blossom round and have the game of forfeits:
        'Spring-happy eyebrow-corners-go-up,' eh?" "That's a fine game of forfeits!"
        Old lady Chia cried, with a smile. "It just suits the time of the year."
        Orders were therefore given at once to fetch a forfeit drum, varnished
        black, and ornamented with designs executed with copper tacks. When brought,
        it was handed to the singing girls to put on the table and rap on it. A twig
        of red plum blossom was then obtained.}
}
{
}

\Verse{73}
{
    \zA{}
}
{
    \eA{"The one in whose hand it is when the drum stops," dowager lady Chia
        laughingly proposed, "will have to drink a cup of wine, and to say something
        or other as well." "I'll tell you what," lady Feng interposed with a smile.
        "Who of us can pit herself against you, dear ancestor, who have ever ready
        at hand whatever you want to say? With the little use we are in this line,
        won't there be an absolute lack of fun in our contributions? My idea is that
        it would be nicer were something said that could be appreciated both by the
        refined as well as the unrefined. So won't it be preferable that the person,
        in whose hands the twig remains, when the drum stops, should crack some joke
        or other?"}
}
{
}

\Verse{74}
{
    \zA{}
}
{
    \eA{Every one, who heard her, was fully aware what a good hand she had always
        been at witty things, and how she, more than any other, had an inexhaustible
        supply of novel and amusing rules of forfeits, ever stocked in her mind, so
        her suggestion not only gratified the various inmates of the family seated
        at the banquet, but even filled the whole posse of servants, both old and
        young, who stood in attendance below, with intense delight. The young
        waiting-maids rushed with eagerness in search of the young ladies and told
        them to come and listen to their lady Secunda, who was on the point again of
        saying funny things. A whole crowd of servant-girls anxiously pressed inside
        and crammed the room. In a little time, the theatricals were brought to a
        close, and the music was stopped.}
}
{
}

\Verse{75}
{
    \zA{}
}
{
    \eA{Dowager lady Chia had some soup, fine cakes and fruits handed to Wen Kuan
        and her companions to regale themselves with, and then gave orders to sound
        the drum. The singing-girls were both experts, so now they beat fast; and
        now slow. Either slow like the dripping of the remnants of water in a
        clepsydra. Or quick, as when beans are being sown. Or with the velocity of
        the pace of a scared horse, or that of the flash of a swift lightning. The
        sound of the drum came to a standstill abruptly. The twig of plum blossom
        had just reached old lady Chia, when by a strange coincidence, the rattle
        ceased. Every one blurted out into a boisterous fit of laughter. Chia Jung
        hastily approached and filled a cup.}
}
{
}

\Verse{76}
{
    \zA{}
}
{
    \eA{"It's only natural," they laughingly cried, "that you venerable senior,
        should be the first to get exhilarated; for then, thanks to you, we shall
        also come in for some measure of good cheer." "To gulp down this wine is an
        easy job," dowager lady smiled, "but to crack jokes is somewhat difficult."
        "Your jokes, dear ancestor, are even wittier than those of lady Feng," the
        party shouted, "so favour us with one, and let's have a laugh!" "I've
        nothing out of the way to evoke laughter with," old lady Chia smilingly
        answered. "Yet all that remains for me to do is to thicken the skin of my
        antiquated phiz and come out with some joke. In a certain family," she
        consequently went on to narrate, "there were ten sons;}
}
{
}

\Verse{77}
{
    \zA{}
}
{
    \eA{these married ten wives. The tenth of these wives was, however, so
        intelligent, sharp, quick of mind, and glib of tongue, that her father and
        mother-in-law loved her best of all, and maintained from morning to night
        that the other nine were not filial. These nine felt much aggrieved and they
        accordingly took counsel together. 'We nine,' they said, 'are filial enough
        at heart; the only thing is that that shrew has the gift of the gab. That's
        why our father and mother-in-law think her so perfect. But to whom can we go
        and confide our grievance?' One of them was struck with an idea. 'Let's go
        to-morrow,' she proposed, 'to the temple of the King of Hell and burn
        incense.}
}
{
}

\Verse{78}
{
    \zA{}
}
{
    \eA{We can then tell the King our grudge and ask him how it was that, when he
        bade us receive life and become human beings, he only conferred a glib
        tongue on that vixen and that we were only allotted such blunt mouths?' The
        eight listened to her plan, and were quite enraptured with it. 'This
        proposal is faultless!' they assented. On the next day, they sped in a body
        to the temple of the God of Hell, and after burning incense, the nine
        sisters-in-law slept under the altar, on which their offerings were laid.
        Their nine spirits waited with the special purpose of seeing the carriage of
        the King of Hell arrive; but they waited and waited, and yet he did not
        come.}
}
{
}

\Verse{79}
{
    \zA{}
}
{
    \eA{They were just giving way to despair when they espied Sun Hsing-che, (the
        god of monkeys), advancing on a rolling cloud. He espied the nine spirits,
        and felt inclined to take a golden rod and beat them. The nine spirits were
        plunged in terror. Hastily they fell on their knees, and pleaded for mercy."
        "'What are you up to?' Sun Hsing-che inquired." "The nine women, with
        alacrity, told him all." "After Sun Hsing-che had listened to their
        confidences, he stamped his foot and heaved a sigh. 'Is that the case?' he
        asked. 'Well, it's lucky enough you came across me, for had you waited for
        the God of Hell, he wouldn't have known anything about it.'" "At these
        assurances, the nine women gave way to entreaties. 'Great saint,' they
        pleaded, 'if you were to display some commiseration, we would be all
        right.'"}
}
{
}

\Verse{80}
{
    \zA{}
}
{
    \eA{"Sun Hsing-che smiled. 'There's no difficulty in the way,' he observed. 'On
        the day on which you ten sisters-in-law came to life, I was, as luck would
        have it, on a visit to the King of Hell's place. So I (saw) him do something
        on the ground, and the junior sister-of-law of yours lap it up. But if you
        now wish to become smart and sharp-tongued, the remedy lies in water. If I
        too were therefore to do something, and you to drink it, the desired effect
        will be attained.'" At the close of her story, the company roared with
        laughter. "Splendid!" shouted lady Feng. "But luckily we're all slow of
        tongue and dull of intellect, otherwise, we too must have had the water of
        monkeys to drink."}
}
{
}

\Verse{81}
{
    \zA{}
}
{
    \eA{"Who among us here," Mrs. Yu and dame Lou smilingly remarked, addressing
        themselves to Li Wan, "has tasted any monkey's water. So don't sham
        ignorance of things!" "A joke must hit the point to be amusing," Mrs. Hsüeh
        ventured. But while she spoke, (the girls) began again to beat the drum. The
        young maids were keen to hear lady Feng's jokes. They therefore explained to
        the singing girls, in a confidential tone, that a cough would be the given
        signal (for them to desist). In no time (the blossom) was handed round on
        both sides. As soon as it came to lady Feng, the young maids purposely gave
        a cough. The singing-girl at once stopped short. "Now we've caught her!"
        shouted the party laughingly; "drink your wine, be quick! And mind you tell
        something nice!}
}
{
}

\Verse{82}
{
    \zA{}
}
{
    \eA{But don't make us laugh so heartily as to get stomachaches." Lady Feng was
        lost in thought. Presently, she began with a smile. "A certain household,"
        she said, "was celebrating the first moon festival. The entire family was
        enjoying the sight of the lanterns, and drinking their wine. In real truth
        unusual excitement prevailed. There were great grandmothers, grandmothers,
        daughters-in-law, grandsons' wives, great grandsons, granddaughters,
        granddaughters-in-law, aunts' granddaughters, cousins' granddaughters; and
        ai-yo-yo, there was verily such a bustle and confusion!" While minding her
        story, they laughed. "Listen to all this mean mouth says!" they cried. "We
        wonder what other ramifications she won't introduce!"}
}
{
}

\Verse{83}
{
    \zA{}
}
{
    \eA{"If you want to bully me," Mrs. Yu smiled, "I'll tear that mouth of yours to
        pieces." Lady Feng rose to her feet and clapped her hands. "One does all one
        can to rack one's brain," she smiled, "and here you combine to do your
        utmost to confuse me! Well, if it is so, I won't go on." "Proceed with your
        story," old lady Chia exclaimed with a smile. "What comes afterwards?" Lady
        Feng thought for a while. "Well, after that," she continued laughingly,
        "they all sat together and crammed the whole room. They primed themselves
        with wine throughout the hours of night and then they broke up."}
}
{
}

\Verse{84}
{
    \zA{}
}
{
    \eA{The various inmates noticed in what a serious and sedate manner she narrated
        her story, and none ventured to pass any further remarks, but waited
        anxiously for her to go on, when they became aware that she coldly and drily
        came to a stop. Shih Hsiang-yün stared at her for ever so long. "I'll tell
        you another," lady Feng laughingly remarked. "At the first moon festival,
        several persons carried a cracker as large as a room and went out of town to
        let it off. Over and above ten thousand persons were attracted, and they
        followed to see the sight. One among them was of an impatient disposition.
        He could not reconcile himself to wait; so stealthily he snatched a
        joss-stick and set fire to it. A sound of 'pu-ch'ih' was heard.}
}
{
}

\Verse{85}
{
    \zA{}
}
{
    \eA{The whole number of spectators laughed boisterously and withdrew. The
        persons, who carried the cracker, felt a grudge against the cracker-seller
        for not having made it tight, (and wondered) how it was that every one had
        left without hearing it go off." "Is it likely that the men themselves
        didn't hear the report?" Hsiang-yün insinuated. "Why, the men themselves
        were deaf," lady Feng rejoined. After listening to her, they pondered for a
        while, and then suddenly they laughed aloud in chorus. But remembering that
        her first story had been left unfinished, they inquired of her: "What was,
        after all, the issue of the first story? You should conclude that too." Lady
        Feng gave a rap on the table with her hand. "How vexatious you are!" she
        exclaimed. "Well, the next day was the sixteenth;}
}
{
}

\Verse{86}
{
    \zA{}
}
{
    \eA{so the festivities of the year were over, and the feast itself was past and
        gone. I see people busy putting things away, and fussing about still, so how
        can I make out what will be the end of it all?" At this, one and all
        indulged in renewed merriment. "The fourth watch has long ago been struck
        outside," lady Feng smilingly said. "From what I can see, our worthy senior
        is also tired out; and we should, like when the cracker was let off in that
        story of the deaf people, be bundling ourselves off and finish!" Mrs. Yu and
        the rest covered their mouths with their handkerchiefs and laughed. Now they
        stooped forward; and now they bent backward. And pointing at her, "This
        thing," they cried, "has really a mean tongue." Old lady Chia laughed.}
}
{
}

\Verse{87}
{
    \zA{}
}
{
    \eA{"Yes," she said, "this vixen Feng has, in real truth, developed a meaner
        tongue than ever! But she alluded to crackers," she added, "so let's also
        let off a few fireworks so as to counteract the fumes of the wine." Chia
        Jung overheard the suggestion. Hurriedly leaving the room, he took the pages
        with him, and having a scaffolding erected in the court, they hung up the
        fireworks, and got everything in perfect readiness. These fireworks were
        articles of tribute, sent from different states, and were, albeit not large
        in size, contrived with extreme ingenuity.}
}
{
}

\Verse{88}
{
    \zA{}
}
{
    \eA{The representations of various kinds of events of antiquity were perfect,
        and in them were inserted all sorts of crackers. Lin Tai-yü was naturally of
        a weak disposition, so she could not stand the report of any loud
        intonation. Her grandmother Chia therefore clasped her immediately in her
        embrace. Mrs. Hsüeh, meanwhile, took Hsiang-yün in her arms. "I'm not
        afraid," smiled Hsiang-yün. "Nothing she likes so much as letting off huge
        crackers," Pao-ch'ai smilingly interposed, "and could she fear this sort of
        thing?" Madame Wang, thereupon, laid hold of Pao-yü, and pulled him in her
        lap. "We've got no one to care a rap for us," lady Feng laughed. "I'm here
        for you," Mrs. Yu rejoined with a laugh. "I'll embrace you. There you're
        again behaving like a spoilt child.}
}
{
}

\Verse{89}
{
    \zA{}
}
{
    \eA{You've heard about crackers, and you comport yourself as if you'd had honey
        to eat! You're quite frivolous again to-day!" "Wait till we break up," lady
        Feng answered laughing, "and we'll go and let some off in our garden. I can
        fire them far better than any of the young lads!" While they bandied words,
        one kind of firework after another was lighted outside, and then later on
        some more again. Among these figured 'fill-heaven-stars;' 'nine
        dragons-enter-clouds;' 'over-whole-land-a- crack-of-thunder;'
        'fly-up-heavens;' 'sound-ten shots,' and other such small crackers. The
        fireworks over, the young actresses were again asked to render the
        'Lotus-flowers-fall,' and cash were strewn upon the stage. The young girls
        bustled all over the boards, snatching cash and capering about. The soup was
        next brought.}
}
{
}

\Verse{90}
{
    \zA{}
}
{
    \eA{"The night is long," old lady Chia said, "and somehow or other I feel
        peckish." "There's some congee," lady Feng promptly remarked, "prepared with
        duck's meat." "I'd rather have plain things," dowager lady Chia answered.
        "There's also some congee made with non-glutinous rice and powder of dates.
        It's been cooked for the ladies who fast." "If there's any of this, it will
        do very well," old lady Chia replied. While she spoke, orders were given to
        remove the remnants of the banquet, and inside as well as outside; were
        served every kind of \_recherché\_ small dishes. One and all then partook of
        some of these refreshments, at their pleasure, and rinsing their mouths with
        tea, they afterwards parted. On the seventeenth, they also repaired, at an
        early hour, to the Ning mansion to present their compliments;}
}
{
}

\Verse{91}
{
    \zA{}
}
{
    \eA{and remaining in attendance, while the doors of the ancestral hall were
        closed and the images put away, they, at length, returned to their quarters.
        Invitations had been issued on this occasion to drink the new year wine at
        Mrs. Hsüeh's residence. But dowager lady Chia had been out on several
        consecutive days, and so tired out did she feel that she withdrew to her
        rooms, after only a short stay. After the eighteenth, relatives and friends
        arrived and made their formal invitations; or else they came as guests to
        the banquets given. But so little was old lady Chia in a fit state to turn
        her mind to anything that the two ladies, Madame Hsing and lady Feng, had to
        attend between them to everything that cropped up.}
}
{
}

\Verse{92}
{
    \zA{}
}
{
    \eA{But Pao-yü as well did not go anywhere else than to Wang Tzu-t'eng's, and
        the excuse he gave out was that his grandmother kept him at home to dispel
        her ennui. We need not, however, dilate on irrelevant details. In due
        course, the festival of the fifteenth of the first moon passed. But, reader,
        if you have any curiosity to learn any subsequent events, listen to those
        given in the chapter below.}
}
{
}
